\documentclass[10pt,letterpaper]{article}
\usepackage{cornell-notes}

\title{Chapter 1: Introduction}
\author{}
\date{}

\setDocTitle{Chapter 1: Introduction}
\setDocSubtitle{Operating Systems}
\setDocVersion{v1.0}
\setDocStatus{Study Companion}
\setDocOwner{Operating Systems Cornell Notes}
\setDocAuthor{}
\setDocDate{\today}
\setDocSummary{Cornell-format study and review notes for Chapter 1: Introduction.}

\setCornellCollection{Operating Systems Cornell Notes}
\setCornellUnitType{Chapter}
\setCornellUnitNumber{1}
\setCornellUnitTitle{Introduction}
\setCornellCompanionLabel{Study and review companion}

\hypersetup{pdftitle={Chapter 1: Introduction},pdfsubject={Cornell notes},pdfauthor={}}

\begin{document}
\maketitle
\section*{Learning Objectives}
\begin{studybox}{By the end of this chapter, you should be able to}
\begin{enumerate}
  \item Explain the operating system as both an extended machine and a resource manager.
  \item Trace the major generations of operating-system development and the hardware changes that enabled them.
  \item Relate processors, memory, disks, I/O devices, buses, and boot firmware to operating-system responsibilities.
  \item Distinguish the operating-system categories described for different workloads and devices.
  \item Explain processes, address spaces, files, I/O, protection, and shells as core abstractions.
  \item Trace representative process, file, directory, and miscellaneous system-call interactions.
  \item Compare monolithic, layered, microkernel, client-server, virtual-machine, and exokernel structures.
  \item Describe why C, headers, compilation units, linking, and runtime layout matter to systems implementation.
\end{enumerate}
\end{studybox}

\begin{studybox}{Section Map}
\hyperlink{sec-1-1}{1.1} \quad \hyperlink{sec-1-2}{1.2} \quad \hyperlink{sec-1-3}{1.3} \quad \hyperlink{sec-1-4}{1.4} \quad \hyperlink{sec-1-5}{1.5} \quad \hyperlink{sec-1-6}{1.6} \quad \hyperlink{sec-1-7}{1.7} \quad \hyperlink{sec-1-8}{1.8} \quad \hyperlink{sec-1-9}{1.9} \quad \hyperlink{sec-1-10}{1.10} \quad \hyperlink{sec-1-11}{1.11} \quad \hyperlink{sec-1-12}{1.12}
\end{studybox}

\section*{Cornell Cue-and-Notes Area}
\small
\begin{longtable}{C N}
\rowcolor{navy}
\color{white}\textbf{Cue / Recall Prompt} &
\color{white}\textbf{Notes}\\
\endfirsthead
\rowcolor{navy}
\color{white}\textbf{Cue / Recall Prompt} &
\color{white}\textbf{Notes (continued)}\\
\endhead
\rowcolor{ice}\multicolumn{2}{l}{\hypertarget{sec-1-1}{}\textbf{Section 1.1: What Is an Operating System?}}\\*\hline
\textbf{1.1.1 Extended machine} & The operating system hides irregular hardware behind stable abstractions. Files, processes, and address spaces let applications use resources without manipulating device registers or raw physical storage.\\\hline
\textbf{1.1.2 Resource manager} & The operating system allocates processors, memory, devices, and storage among competing programs. It multiplexes resources in time, such as CPU turns, and in space, such as memory regions or disk blocks.\\\hline
\textbf{Two complementary views} & The abstraction view explains the interface presented upward to programs; the resource-management view explains control of hardware below. A practical system must satisfy both views at once.\\\hline
\rowcolor{ice}\multicolumn{2}{l}{\hypertarget{sec-1-2}{}\textbf{Section 1.2: History of Operating Systems}}\\*\hline
\textbf{1.2.1 Vacuum-tube era} & \textcolor{legacy}{\textbf{Historical or legacy context:}} Early machines had no operating system. A single team wired or entered one machine-language job at a time, so setup, operation, and computation were inseparable.\\\hline
\textbf{1.2.2 Batch systems} & \textcolor{legacy}{\textbf{Historical or legacy context:}} Transistors improved reliability and enabled professional operators, card input, magnetic-tape staging, and monitors that automatically sequenced batches of jobs.\\\hline
\textbf{1.2.3 ICs and multiprogramming} & \textcolor{legacy}{\textbf{Historical or legacy context:}} Integrated circuits supported product families and multiprogramming. Spooling, memory protection, interrupts, and time sharing improved utilization while allowing interactive access.\\\hline
\textbf{1.2.4 Personal computers} & \textcolor{legacy}{\textbf{Historical or legacy context:}} Microprocessors moved computing to individual users. Command-line systems evolved toward graphical interfaces, while networked and UNIX-derived systems spread to workstations and servers.\\\hline
\textbf{1.2.5 Mobile computers} & \textcolor{legacy}{\textbf{Historical or legacy context:}} Mobile systems combined constrained energy budgets, wireless connectivity, sensors, and touch interfaces, requiring different application and resource-management priorities.\\\hline
\rowcolor{ice}\multicolumn{2}{l}{\hypertarget{sec-1-3}{}\textbf{Section 1.3: Computer Hardware Review}}\\*\hline
\textbf{1.3.1 Processors} & A CPU fetches, decodes, and executes instructions; pipelines, superscalar execution, privilege modes, traps, and multicore designs affect how the kernel obtains control and protects itself.\\\hline
\textbf{1.3.2 Memory} & Registers, caches, main memory, and secondary storage form a hierarchy: faster levels are smaller and costlier. The memory-management unit translates addresses and enforces access restrictions.\\\hline
\textbf{1.3.3 Disks} & Magnetic disks organize data into sectors and tracks and incur seek and rotational delays; solid-state storage has different timing and wear properties. Controllers expose logical operations to the operating system.\\\hline
\textbf{1.3.4 I/O devices} & Controllers and device registers mediate I/O. Programmed I/O, interrupts, and direct memory access trade CPU involvement against controller complexity and transfer efficiency.\\\hline
\textbf{1.3.5 Buses} & Buses and point-to-point interconnects carry addresses, data, and control signals among processors, memory, and devices. Arbitration and bandwidth constrain overall system throughput.\\\hline
\textbf{1.3.6 Booting} & Firmware begins execution, discovers or initializes hardware, locates a bootable device, loads a bootstrap program, and ultimately transfers control to the operating-system kernel.\\\hline
\rowcolor{ice}\multicolumn{2}{l}{\hypertarget{sec-1-4}{}\textbf{Section 1.4: The Operating System Zoo}}\\*\hline
\textbf{1.4.1 Mainframes} & Mainframe systems emphasize high I/O throughput, batch processing, transaction workloads, reliability, and extensive resource accounting.\\\hline
\textbf{1.4.2 Servers} & Server systems provide shared services to many clients and emphasize concurrency, availability, networking, protection, and scalable storage.\\\hline
\textbf{1.4.3 Multiprocessors} & Multiprocessor systems coordinate multiple CPUs sharing resources. Scheduling, cache coherence, synchronization, and scalable kernel data structures become central.\\\hline
\textbf{1.4.4 Personal computers} & Personal-computer systems balance interactive responsiveness, device support, usability, networking, and application compatibility for one primary user at a time.\\\hline
\textbf{1.4.5 Handhelds} & Handheld systems manage batteries, radios, sensors, touch input, limited displays, and application life cycles under tight energy and memory constraints.\\\hline
\textbf{1.4.6 Embedded systems} & Embedded systems are built into products, often execute a fixed function, and may have constrained memory, specialized devices, or strict reliability requirements.\\\hline
\textbf{1.4.7 Sensor nodes} & Sensor-node systems coordinate small, networked devices with severe energy and storage limits; communication and duty cycling can dominate resource policy.\\\hline
\textbf{1.4.8 Real-time systems} & Real-time correctness includes timing. Hard real-time systems must meet deadlines; soft real-time systems tolerate occasional misses with degraded value rather than total failure.\\\hline
\textbf{1.4.9 Smart cards} & Smart-card systems run with tiny memory and processing budgets and may host a restricted interpreter or several protected applets.\\\hline
\rowcolor{ice}\multicolumn{2}{l}{\hypertarget{sec-1-5}{}\textbf{Section 1.5: Operating System Concepts}}\\*\hline
\textbf{1.5.1 Processes} & A process is a program in execution plus its address space, registers, open resources, and scheduling state. Creation and termination form dynamic process populations.\\\hline
\textbf{1.5.2 Address spaces} & An address space is the set of memory locations a process may address. Isolation prevents unrelated processes from corrupting one another while controlled sharing enables cooperation.\\\hline
\textbf{1.5.3 Files} & Files provide persistent named byte or record collections. Directories organize names hierarchically; links, permissions, descriptors, and mounting shape access.\\\hline
\textbf{1.5.4 Input/output} & I/O abstractions and drivers isolate applications from device-specific protocols. The kernel coordinates sharing, buffering, naming, errors, and completion notification.\\\hline
\textbf{1.5.5 Protection} & Protection defines which subjects may perform which operations on which objects. Privilege modes, authentication, authorization, and isolation enforce boundaries.\\\hline
\textbf{1.5.6 Shell} & A shell is a user-level command interpreter, not the kernel. It starts programs, redirects streams, constructs pipelines, and provides scripting control structures.\\\hline
\textbf{1.5.7 Recurring evolution} & \textcolor{legacy}{\textbf{Historical or legacy context:}} Constraints can cause old architectural ideas to reappear at new scales; layered interpreters, compact systems, and centralized services recur as hardware contexts change.\\\hline
\rowcolor{ice}\multicolumn{2}{l}{\hypertarget{sec-1-6}{}\textbf{Section 1.6: System Calls}}\\*\hline
\textbf{System-call path} & A library wrapper places arguments where the calling convention requires, executes a trap, and lets the kernel validate and dispatch the request before returning a result or error.\\\hline
\textbf{1.6.1 Process calls} & Representative operations create a child, replace a process image, wait for completion, obtain identifiers, and terminate. \texttt{fork}, \texttt{exec}, \texttt{waitpid}, and \texttt{exit} illustrate separation of creation from program loading.\\\hline
\textbf{1.6.2 File calls} & Opening associates a path with a descriptor; reading and writing transfer bytes; seeking changes the current offset; metadata operations inspect or modify attributes; closing releases the descriptor.\\\hline
\textbf{1.6.3 Directory calls} & Directory operations create, remove, link, unlink, rename, and change the process working directory. Path resolution walks directory components subject to permissions.\\\hline
\textbf{1.6.4 Miscellaneous calls} & Other calls obtain time or system information, change protection attributes, manage signals, or control devices when a uniform read/write interface is insufficient.\\\hline
\textbf{1.6.5 Win32 API} & \textcolor{legacy}{\textbf{Historical or legacy context:}} The Win32 interface is a broad library API rather than a one-to-one list of kernel calls; some procedures are implemented in user space and others invoke native services.\\\hline
\rowcolor{ice}\multicolumn{2}{l}{\hypertarget{sec-1-7}{}\textbf{Section 1.7: Operating System Structure}}\\*\hline
\textbf{1.7.1 Monolithic systems} & Most services execute in one kernel address space and call one another directly. This offers efficient communication but makes fault containment and dependency control difficult.\\\hline
\textbf{1.7.2 Layered systems} & Each layer uses services below and exports services upward. Clear interfaces improve reasoning, although real dependencies do not always fit a strict hierarchy.\\\hline
\textbf{1.7.3 Microkernels} & A microkernel retains minimal mechanisms such as address spaces, scheduling, and interprocess communication while moving servers to user mode. Isolation improves reliability at a communication cost.\\\hline
\textbf{1.7.4 Client-server model} & Clients request services from separate server processes through messages. The model supports modularity and distribution when naming, failure, and communication semantics are defined carefully.\\\hline
\textbf{1.7.5 Virtual machines} & A virtualization layer presents multiple isolated machine interfaces on one physical computer. Each guest can run its own operating system while the monitor controls resources.\\\hline
\textbf{1.7.6 Exokernels} & An exokernel securely multiplexes hardware while delegating higher-level abstractions to library operating systems, giving applications more control at the cost of a different programming model.\\\hline
\rowcolor{ice}\multicolumn{2}{l}{\hypertarget{sec-1-8}{}\textbf{Section 1.8: The World According to C}}\\*\hline
\textbf{1.8.1 C language} & C exposes addresses, memory layout, bit operations, and low-level control with limited runtime protection, making it suitable for kernels but demanding disciplined handling of types and storage.\\\hline
\textbf{1.8.2 Header files} & Headers declare types, constants, macros, and external interfaces shared across translation units. Include guards and consistent declarations prevent conflicting views of an interface.\\\hline
\textbf{1.8.3 Large projects} & Systems code is divided into source files, compiled to object files, and linked with libraries. Build dependencies, configuration variants, and modular interfaces control complexity.\\\hline
\textbf{1.8.4 Runtime model} & A process image typically contains text, initialized data, uninitialized data, heap, and stack regions. Calls create stack frames; dynamic allocation grows managed storage independently.\\\hline
\rowcolor{ice}\multicolumn{2}{l}{\hypertarget{sec-1-9}{}\textbf{Section 1.9: Research on Operating Systems}}\\*\hline
\textbf{Research directions} & Research tests new structures, reliability techniques, security mechanisms, scheduling policies, virtualization methods, and support for emerging hardware. Evaluation must connect a proposed mechanism to a measurable goal.\\\hline
\rowcolor{ice}\multicolumn{2}{l}{\hypertarget{sec-1-10}{}\textbf{Section 1.10: Outline of the Remaining Material}}\\*\hline
\textbf{Conceptual roadmap} & The subsequent sequence develops processes and threads, memory, file systems, I/O, deadlocks, virtualization, multiprocessors and distributed systems, security, concrete system case studies, and design principles.\\\hline
\rowcolor{ice}\multicolumn{2}{l}{\hypertarget{sec-1-11}{}\textbf{Section 1.11: Metric Units}}\\*\hline
\textbf{Powers and prefixes} & Decimal prefixes normally describe rates and some device capacities, while binary powers often describe memory. Context matters: a nominal prefix may not imply the same base in every hardware discussion.\\\hline
\rowcolor{ice}\multicolumn{2}{l}{\hypertarget{sec-1-12}{}\textbf{Section 1.12: Summary}}\\*\hline
\textbf{Foundation} & Operating systems provide abstractions, manage resources, mediate protected transitions through system calls, and organize implementation using structures chosen for performance, reliability, and maintainability.\\\hline
\end{longtable}
\normalsize

\section*{Chapter Synthesis}
\begin{studybox}{Chapter Summary}
An operating system simultaneously makes hardware easier to use and arbitrates access to it. Its evolution follows changes in hardware and workloads, but recurring abstractions--processes, address spaces, files, I/O, and protection--remain central. System calls define the application-visible contract, while kernel structure determines how responsibilities are divided internally.
\end{studybox}

\begin{studybox}{Key Relationships}
\begin{itemize}
  \item Hardware privilege and interrupts make protected system-call entry and resource arbitration possible.
  \item Processes use address spaces, files, and I/O abstractions; protection constrains every relationship.
  \item System-call interfaces expose services whose implementations depend on the chosen kernel structure.
  \item Historical workload changes explain why different operating-system categories emphasize different policies.
\end{itemize}
\end{studybox}

\begin{studybox}{Design Tradeoffs}
\begin{itemize}
  \item Rich abstractions improve portability and usability but can hide costs or restrict specialized control.
  \item Monolithic execution reduces communication overhead, whereas smaller protected components improve isolation.
  \item Resource utilization can conflict with responsiveness, predictability, energy use, or fairness.
  \item Compatibility preserves software investment but constrains redesign and simplification.
\end{itemize}
\end{studybox}

\begin{studybox}{Common Confusions}
\begin{itemize}
  \item A shell is an ordinary user program that uses system calls; it is not the operating system kernel.
  \item An API procedure is not necessarily a system call, and one API call may require zero, one, or several kernel entries.
  \item Multiprogramming improves utilization; time sharing adds rapid switching to support interactive users.
\end{itemize}
\end{studybox}

\section*{Key Terms}
\small
\begin{longtable}{C N}
\rowcolor{navy}\color{white}\textbf{Term} & \color{white}\textbf{Concise definition}\\
\endfirsthead
\rowcolor{navy}\color{white}\textbf{Term} & \color{white}\textbf{Concise definition (continued)}\\
\endhead
\textbf{Kernel} & Privileged core that manages hardware and implements protected operating-system services.\\\hline
\textbf{Process} & Program in execution together with its state and owned resources.\\\hline
\textbf{Address space} & Memory addresses a process is permitted to reference.\\\hline
\textbf{System call} & Controlled request from a program to the kernel.\\\hline
\textbf{Trap} & Synchronous control transfer into privileged software.\\\hline
\textbf{Interrupt} & Asynchronous hardware event that transfers control to an interrupt handler.\\\hline
\textbf{Multiprogramming} & Keeping multiple jobs available so another can run when one waits.\\\hline
\textbf{Time sharing} & Rapid processor sharing intended to provide interactive service.\\\hline
\textbf{Device driver} & Software that controls a device and exposes an operating-system interface.\\\hline
\textbf{DMA} & Controller-driven block transfer between a device and memory.\\\hline
\textbf{Microkernel} & Minimal privileged kernel with additional services in isolated processes.\\\hline
\textbf{Virtual machine} & Software-defined machine interface isolated from other guests.\\\hline
\textbf{Shell} & User-level command interpreter and scripting environment.\\\hline
\textbf{Bootstrapping} & Staged process that initializes hardware and loads the kernel.\\\hline
\end{longtable}
\normalsize

\enlargethispage{2\baselineskip}
\begin{studybox}{Active-Recall Review}
\small
\begin{enumerate}
  \item How do the extended-machine and resource-manager views complement one another?
  \item Why did interrupts, protection, and multiprogramming develop together?
  \item What steps transfer control from firmware to a running kernel?
  \item How do hard and soft real-time requirements differ?
  \item What state must accompany a program for it to be treated as a process?
  \item Why is an API not identical to a system-call interface?
  \item How do monolithic and microkernel structures trade communication cost for fault isolation?
  \item Why does systems programming require explicit attention to runtime memory layout?
\end{enumerate}
\end{studybox}

\par\medskip
{\color{navy}\bfseries Next-Chapter Connections}\par\smallskip
Chapter 2 turns the introductory process abstraction into a concrete model of process states, threads, communication, synchronization, and scheduling.
\normalsize
\end{document}
